\chapter{Verbs with Objects: Akkusativ, Dativ, and Two-Way}
\label{cha:VerbenObjekte}

This chapter provides comprehensive verb tables with their required cases and example sentences - essential for German grammar.

\section{Verbs with Akkusativ (Accusative Verbs)}
\label{sec:AkkusativVerben}

These verbs require an Akkusativ object (Wen? Was?):

\bigskip
\noindent\textbf{Common Akkusativ Verbs:}

\begin{center}
\begin{tabular}{|l|c|l|}
\hline
Verb & Example & Meaning \\ \hline
sehen & Ich sehe den Mann. & see \\ \hline
hören & Ich höre die Musik. & hear \\ \hline
lesen & Er liest das Buch. & read \\ \hline
schreiben & Sie schreibt einen Brief. & write \\ \hline
suchen & Ich suche den Schlüssel. & search for \\ \hline
finden & Wir finden die Lösung. & find \\ \hline
kennen & Ich kenne sie. & know (person) \\ \hline
treffen & Ich treffe ihn morgen. & meet \\ \hline
brauchen & Ich brauche das Buch. & need \\ \hline
haben & Er hat ein Auto. & have \\ \hline
kaufen & Sie kauft die Blumen. & buy \\ \hline
essen & Ich esse den Apfel. & eat \\ \hline
trinken & Er trinkt den Kaffee. & drink \\ \hline
verstehen & Ich verstehe die Frage. & understand \\ \hline
lieben & Ich liebe dich. & love \\ \hline
mögen & Ich mag das Essen. & like \\ \hline
\end{tabular}
\end{center}

\bigskip
\noindent\textbf{With Example Sentences:}

\begin{enumerate}
    \item \textit{Ich sehe meinen Freund jeden Tag.} (I see my friend every day.)
    \item \textit{Sie liest die Zeitung am Morgen.} (She reads the newspaper in the morning.)
    \item \textit{Wir brauchen mehr Zeit.} (We need more time.)
    \item \textit{Er hat ein neues Auto gekauft.} (He bought a new car.)
    \item \textit{Ich verstehe die Aufgabe nicht.} (I don't understand the task.)
\end{enumerate}

\section{Verbs with Dativ (Dative Verbs)}
\label{sec:DativVerben}

These verbs require a Dativ object (Wem?):

\bigskip
\noindent\textbf{Common Dativ Verbs:}

\begin{center}
\begin{tabular}{|l|c|l|}
\hline
Verb & Example & Meaning \\ \hline
helfen & Ich helfe der Frau. & help \\ \hline
danken & Ich danke dir. & thank \\ \hline
gefallen & Das Buch gefällt mir. & be pleasing \\ \hline
gehören & Das Auto gehört ihm. & belong to \\ \hline
fehlen & Mir fehlt die Zeit. & be missing \\ \hline
passen & Die Schuhe passen mir. & fit \\ \hline
schmecken & Der Kuchen schmeckt uns. & taste \\ \hline
glauben & Ich glaube ihm. & believe \\ \hline
zustimmen & Ich stimme dir zu. & agree \\ \hline
folgen & Der Hund folgt ihr. & follow \\ \hline
antworten & Er antwortet dem Lehrer. & answer \\ \hline
gratulieren & Ich gratuliere ihr. & congratulate \\ \hline
zuhören & Du hörst mir zu. & listen to \\ \hline
begegnen & Ich begegne ihr in der Stadt. & encounter \\ \hline
\end{tabular}
\end{center}

\bigskip
\noindent\textbf{With Example Sentences:}

\begin{enumerate}
    \item \textit{Ich helfe meiner Mutter bei der Arbeit.} (I help my mother with the work.)
    \item \textit{Vielen Dank für Ihre Hilfe!} (Thank you very much for your help!)
    \item \textit{Das neue Kleid gefällt mir sehr.} (I really like the new dress.)
    \item \textit{Ich stimme Ihrer Meinung zu.} (I agree with your opinion.)
    \item \textit{Entschuldigung, ich höre Ihnen nicht zu.} (Sorry, I'm not listening to you.)
\end{enumerate}

\section{Verbs with Dativ and Akkusativ (Two-Way Verbs)}
\label{sec:ZweiwegVerben}

These verbs require both a Dativ and an Akkusativ object:

\bigskip
\noindent\textbf{Common Two-Way Verbs:}

\begin{center}
\begin{tabular}{|l|c|c|l|}
\hline
Verb & Dativ & Akkusativ & Meaning \\ \hline
geben & mir & das Buch & give \\ \hline
bringen & ihr & die Nachricht & bring \\ \hline
sagen & mir & die Wahrheit & tell \\ \hline
schicken & mir & eine E-Mail & send \\ \hline
zeigen & uns & den Weg & show \\ \hline
erklären & dem Lehrer & die Aufgabe & explain \\ \hline
kaufen & mir & ein Geschenk & buy \\ \hline
schenken & ihr & Blumen & give as gift \\ \hline
empfehlen & mir & ein Buch & recommend \\ \hline
erzählen & mir & eine Geschichte & tell (story) \\ \hline
\end{tabular}
\end{center}

\bigskip
\noindent\textbf{With Example Sentences:}

\begin{enumerate}
    \item \textit{Ich gebe dem Kind ein Buch.} (I give the child a book.)
    \item \textit{Sie bringt ihrer Mutter Blumen.} (She brings her mother flowers.)
    \item \textit{Er erklärt dem Studenten die Lösung.} (He explains the solution to the student.)
    \item \textit{Ich schenke meiner Freundin ein Schmuck.} (I give my girlfriend jewelry as a gift.)
    \item \textit{Kannst du mir bitte den Weg zeigen?} (Can you please show me the way?)
\end{enumerate}

\section{Verbs with Genitiv}
\label{sec:GenitivVerben}

These verbs require a Genitiv object (less common, more formal):

\bigskip
\begin{center}
\begin{tabular}{|l|c|l|}
\hline
Verb & Example & Meaning \\ \hline
bedürfen & Das bedarf keiner Erklärung. & require \\ \hline
gedenken & Wir gedenken der Opfer. & commemorate \\ \hline
sich bedienen & Er bedient sich der Hilfe. & make use of \\ \hline
sich erfreuen & Das Projekt erfreut sich großer Beliebtheit. & enjoy \\ \hline
sich annehmen & Sie nimmt sich des Problems an. & take care of \\ \hline
\end{tabular}
\end{center}

\chapterend